\chapter{Компьютер из живых клеток}
% полная версия: chapters/21_living.tex

Можно ли построить машину не из кремния, а из живых нейронов? Можно, и её уже строят. Нейроны выращивают прямо на чипе с тысячами электродов --- плоским слоем или комочками, похожими на крошечный мозг. Каждый электрод умеет и слушать, и подавать сигнал. Это стенд с открытой схемой: всё можно измерить и во всё можно вмешаться. Только радиостанций в нём нет --- это шина данных без регистров управления.

\section*{Что сеть делает сама}

Предоставленная самой себе, культура время от времени вспыхивает целиком: почти все нейроны щёлкают вместе, а потом затихают, --- с промежутками от секунды до минут. Это уже знакомый генератор: сеть возбуждает сама себя и устаёт, пока соединения не восстановят запас вещества.

\section*{Учитель без дофамина}

Как научить такую сеть, если дофамина нет? Первые способы были электрическими. Например: стимулировать сеть, пока она не выдаст нужный ответ, --- и сразу остановиться. Ответ закрепляется. В другом опыте культура управляла простым виртуальным «телом» и училась вести его к цели. Во всех этих опытах учитель --- рисунок тока, который выбирает инженер.

\section*{Дофамин по вспышке света}

Однажды учителя попробовали сделать химическим. На одной из платформ с живыми комочками нейронов в питательную жидкость добавили дофамин в светочувствительной «клетке», которая не даёт ему действовать. Когда сеть отвечала на стимул сильнее прежнего, её «награждали»: вспышка ультрафиолета разрушала клетку, и дофамин выходил наружу. За двести сорок повторов ответы в целом росли. Но сами авторы честно пишут, что по одному такому опыту нельзя сказать, дофамин ли тому причиной.

\section*{Шест на тележке}

В недавнем опыте комочки нейронов учили балансировать шестом на тележке --- классическая задача для программ обучения. Какие учебные импульсы подавать, выбирала программа. Выбранные программой сигналы работали лучше случайных, но через три четверти часа отдыха выученное исчезало. А если заблокировать \nt{ГЛУТ}-соединения, обучение не шло вовсе: выученное живёт именно там. Учитель же по-прежнему снаружи --- только теперь это не инженер, а программа.

\section*{Живой резервуар}

Есть и другой подход: живую сеть не учат вовсе. Её используют как «резервуар» --- богатую, запутанную систему, чьи ответы простой электронный счётчик учится разбирать. Так живой комочек помогал различать голоса. Правда, по словам авторов, связи в самом комочке во время работы тоже менялись.

\section*{Цена и вопрос}

Миллион нейронов на чипе потребляет примерно десятитысячную долю ватта --- куда меньше, чем электроника вокруг них. И стенд не умеет сам решить, что считать успехом: эту роль всегда играет кто-то снаружи. А по мере того как такие системы растут, встаёт и вопрос этики, который пока никто не решил.

\fullversion{21}
